\documentclass[a4paper,12pt]{jarticle}

\topmargin=0mm \voffset = -5mm \headheight=5mm \headsep = 10mm 
\textheight = 235mm \textwidth = 160mm \topskip = 3mm 
\evensidemargin = 0mm \oddsidemargin = 0mm \footskip = 15mm 

\usepackage[dvipdfmx]{graphicx}
\usepackage {amsmath,amssymb,multicol,amsthm}
%,epic,eepic,ecltree \usepackage{listings} めんどくさいんでverbatim
\special{pdf: minorversion=6}%自宅での対応
\usepackage{fancyhdr}
\usepackage{here}
\pagestyle{fancy}
\rhead{浜田宏}
\lhead{　The One with...}
%\rhead{2014メモ　浜田宏}


\theoremstyle{definition}%定理型の環境を全て立体指定にする
\newtheorem{definition}{定義}
\newtheorem{prop}{命題}
\newtheorem{Axiom}{公理}
\newtheorem{ass}{仮定}
\newtheorem{col}[prop]{系}
\newtheorem{EX}{例}
\newtheorem{note}{注意}
\newtheorem{thm}[prop]{定理}
\newtheorem{exercise}{問題}
\newtheorem{answer}{解答例}
\newtheorem{lemma}{補題}


\newtheorem{pro}{命題}%番号無しの命題環境
\renewcommand{\thepro}{}

\newtheorem{colo}{系}%番号無しの系環境
\renewcommand{\thecolo}{}

\newtheorem{defi}{定義}%番号無しの定義
\renewcommand{\thedefi}{}

\newcommand{\bib}{\hangafter=1\hangindent=11mm\noindent}
%本文では，つぎのように使う．\section*{文献}
%\bib Axelrod, Robert. 1984. {\it The Evolution of Cooperation}. Basic Books.

\newcommand{\F}{\mathcal{F}}%加算加法族
\newcommand{\BO}{\mathcal{B}}%ボレル集合体
\newcommand{\Om}{\Omega}
\newcommand{\w}{\omega}
\newcommand{\mcc}{\mathcal}%mathカリグラフィック

\newcommand{\e}{\varepsilon}
\newcommand{\ra}{\rightarrow}
\newcommand{\lra}{\Longrightarrow}

\newcommand{\mvec}{\boldsymbol}
\newcommand{\gvec}[1]{\ensuremath{{\boldsymbol{#1}=(#1_{1},#1_{2}, \ldots, #1_{n})}}}
\newcommand{\D}{\displaystyle}
\newcommand{\up}{\upshape}
\newcommand{\bp}{\begin{proof}[\textbf{証明}]\upshape}
\newcommand{\ep}{\end{proof}}
\newcommand{\set}[1]{\{#1\}}
\newcommand{\sset}[3]{\{#1,#2,\ldots,#3\}}
\newcommand{\bd}{\begin{definition}\upshape}
\newcommand{\ed}{\end{definition}}
\newcommand{\mcal}{\mathcal}
\newcommand{\B}{\textbf}
\newcommand{\GG}{\Gamma}

\newenvironment{nitem}{\begin{list}
{\textbullet}
{\setlength{\itemsep}{0cm} \setlength{\parsep}{0cm}}}{\end{list}}

\newcommand{\nele}[1]{\ensuremath{#1_{1},#1_{2}, \ldots, #1_{n}}}
\newcommand{\nvec}[1]{\ensuremath{(#1_{1},#1_{2}, \ldots, #1_{n})}}
\newcommand{\nset}[1]{\ensuremath{ \{ #1_{1},#1_{2}, \ldots, #1_{n} \}}}
\newcommand{\A}{\forall}
\newcommand{\E}{\exists}
\newcommand{\q}{\hspace*{1em}} %ちょいスペース用
\renewcommand{\baselinestretch}{1.25} %行間調整


\let\oldenumerate\enumerate
\renewcommand{\enumerate}{
	\oldenumerate
	\setlength{\itemsep}{1pt}
	\setlength{\parskip}{0pt}
	\setlength{\parsep}{0pt}
}
\let\olditemize\itemize
\renewcommand{\itemize}{
	\olditemize
	\setlength{\itemsep}{1pt}
	\setlength{\parskip}{0pt}
	\setlength{\parsep}{0pt}
}

%\title{朝食付きローマと朝食無しローマ\\Toy model series 1}
%\author{浜田　宏\\ 東北大学}
%\date{}
%%%%%%    TEXT START    %%%%%%
\begin{document}



\section{動機}
	
%	拳銃を真上に向かって撃ち，落ちてきた弾に当たるとどうなるのか？}


子供の頃からずっと疑問に思っていることがあった．


それは

\begin{quote}「拳銃を真上に向かって撃ち，落ちてきた弾に当たるとどうなるのか？」
\end{quote}

という問題である．

\bigskip

落ちてきた弾のスピードはどのくらいだろうか？

弾が落ちてくるまでどのくらい時間がかかるのだろうか？

落ちてきた弾に当たると痛いのだろうか？

それとも別段大したことはないのだろうか？

$\vdots$

疑問は尽きない．

$\vdots$


そしてこの疑問が解けないとたいへん困る．

このままだとゾンビに取り囲まれたとき，威嚇射撃として頭上に拳銃を撃つことができないではないか．

たいへんだ．

どうやって血に飢えたゾンビから我が身を守ればよいのだ．

%と私は心に決めている．
%ところが撃った後の弾の行方が気になって，いざというとき頭上に向かって撃てないと困るではないか．

というわけで計算して確かめることにした．簡単のため，まずは空気抵抗はないものと考える．



\section{ニュートンの運動方程式}

ニュートンの運動方程式をつかって計算すると，真上に向かって初速$v_0$で投げた物体の$t$秒後の位置は
\[ 
y(t)=-\frac{1}{2}gt^2+v_0 t
\]
である（$g$は重力加速度）．

%\[ 
%y(t)=-\frac{1}{2}(9.8m/s^2)t^2+v_0 t
%\]
%である．

拳銃で発射した弾丸の速度を調べたところ，概ね$360m/s \sim 420m/s$の範囲で分布しているらしい．そこで計算しやすさを考えて弾丸の初速$v_0$を$v_0=400m/s$と仮定する．

真上に向かって発射した弾丸は下向きの力（重力）によってだんだんと速度が低下して，0に近づき，やがて落下する．
$y=0$が，地上に落下した位置なので，そこに到達する時間は
\[ 
-\frac{1}{2}g t^2+ (400m/s)t=0
\]
を満たす時刻$t$である．$t=0$は自明なので，この時刻を除いて計算すると
\begin{align*} 
%-t(\frac{1}{2}(9.8m/s^2)t-400m/s)& = 0\\ 
\frac{1}{2}gt-400m/s& = 0\\
t&=\frac{800m/s}{g}
\end{align*}
重力加速度を$g=9.8m/s^2$とおけば
\[ 
t=\frac{800m/s}{9.8m/s^2}=\frac{400}{4.9}s=81.6327s
\]
である．

つまり$81.6327$秒後に地上に落ちてくる，というわけだ．

これだけ時間があればとりあえず逃げる時間は十分にありそうだ．


よかった．


これでゾンビに取り囲まれたとき，安心して真上に向かって威嚇射撃ができる．


\section{弾はどこまで上昇するのか？}


さて，真上に向かって初速$400m/s$で弾丸を発射すると，$81.6327$秒後に地上に落ちてくることが分かった．

落ちてくる，ということは，どこかで最高高度に達するはずだ．

一体その高度はどのくらいであろうか？

たとえば私が出張中に飛行機に乗っているとき，その真下にゾンビに取り囲まれた人がいたとする．その人が威嚇射撃として真上に銃を撃ったとき，飛行機に乗った私にその弾が当たるのだろうか？

当たったらどうしよう．

心配で，出張用に持ってきた専門書を機内でゆっくり読めないではないか．

というわけで弾丸の最高高度を計算してみよう．

時間$t$における弾丸の位置$y$は
\[ 
y(t)=-\frac{1}{2}gt^2+v_0 t
\]
であった．この$y(t)$を時間$t$で微分すれば，速度$v(t)$を得る．

すなわち速度$v(t)$は
\[ 
v(t)=\frac{dy}{dt}=-gt+v_0
\]
である．

最高高度に到達したとき速度$v(t)$は0になるので
\[ 
-gt+v_0=0 \iff t=\frac{v_0}{g} 
\]
ここに初速度$v_0$と重力加速度$g=9.8m/s^2$を代入すれば
\[ 
t=\frac{400m/s}{9.8m/s^2}=40.8163s
\]
となる．つまり弾丸を発射してから40.8163秒後に最高高度に到達するはずである．

このとき弾丸の位置$y(t)$は
\[ 
y(t)=-\frac{1}{2}g(40.8163s)^2+v_0(40.8163s)=8163.27m
\]
となる．

想像してほしい．

真上に向かって拳銃をバァーンと撃つ．すると弾丸は初速$400m/s$で上昇するが，重力によってだんだんと速さが遅くなる．そして約40秒後には速度が0になり，ぴたりと止まる．

そのときの弾丸の高度は地上から$8163.27m$である\footnote{もちろん実際には空気抵抗などがあるから，そこまでは飛ばないだろう．これについては5節で考える}．

ちなみに航空機の高度は，およそ$8000m$から$12000m$なので，地上から飛行機を狙って拳銃で撃つと空気抵抗がなければ理論上は弾が当たると言える．

もっともその高度で当たったとしても弾丸の速度が重力によって0近くに減速しているので，航空機なら跳ね返すだろう．

よかった．これで安心して機内で『HUNTER $\times$ HUNTER』を読むことができる．



\section{落ちてきた弾に当たるとどうなるのか？}


空気抵抗がないとき真上に向かって撃った弾丸は$8163.27m$まで上昇し，やがて地上に落ちてくることが分かった．

次に，地上に落ちてきた弾にどのくらい威力があるのかを考えてみよう．

先ほどの計算から，弾丸が地上に到達する時間は，$81.6327$秒後である．

この時の弾丸の速度は
\[ 
v(t)=-gt+v_0=-(9.8m/s^2)81.6327s+400m/s=-800m/s+400m/s=-400m/s
\]
である．

速度にマイナスがついているのは，上向きを正と考えていたからで，速度$-400m/s$は下向きに初速と同じ速さで落ちてくる，という意味である．

おもしろい．


初速$400m/s$で撃った弾が，きっちり同じ速さで地上に戻ってくるのだ．

このことは偶然ではない．


一般に，発射した弾が地上に戻ってくる時間は$y(t)=0$を満たす$t$である．つまり
\begin{align*}
y(t)&=0  \\ 
-\frac{1}{2}gt^2+v_0 t&=0\\
-t(\frac{1}{2}gt-v_0)&=0\\
t&=\frac{2v_0}{g}   
\end{align*} 
に地上に戻ってくる．このときの速度は
\[ 
v(t)=-gt+v_0=-g\frac{2v_0}{g}+v_0=-v_0
 \]
である．つまり必ず反対方向（下向き）の同じ速さとなる．

なお速度0に到達する時間が40.8163秒で，地上に落ちてくる時間が81.6327秒だから，ちょうど
\[ 
40.8163秒 \times 2=81.6327秒
\]
となっており，最高高度に達する時間とそこから地上に戻ってくる時間が，ぴったり一致している．

このことも偶然ではない．

というのも一般に，最高高度に到達する時点で速度は0になるから
\begin{align*}
v(t)&=0   \\ 
-gt+v_0&=0\\
t=\frac{v_0}{g}   
\end{align*} 
より，最高高度に達するまでの時間は$\frac{v_0}{g}$である．

一方地上に戻ってくる時間は
\[ t=\frac{2v_0}{g}   
 \]
である．

よって「最高高度に達するまでの時間」の2倍が，「地上に戻ってくる時間」に一致する．

言い換えれば「最高高度に達するまでの時間」と「最高高度から自由落下で地上に戻ってくるまでの時間」は一致する．

まったく，うまくできている．

このことを高校の物理の先生がもっと分かりやすく教えてくれていたら，私は理系学部に進学していたかもしれない．


さて，地上に戻ってきた弾丸の速度が初速と同じであること，は何を意味するだろうか？

運動エネルギー$K$は物体の質量$m$に比例し，速さ$v$の2乗に比例する．
\[ 
K=\frac{1}{2}m v^2
\]
ということは$\cdots \cdots$．

落ちてきた弾の運動エネルギーは，発射した瞬間の弾丸の運動エネルギーと同じではないか．

なんということだ．

落ちてきた弾に当たることは，撃った瞬間の弾に当たるのと同じ運動エネルギーを受けることを意味する．

これは相当に危険だ．

これからはゾンビに威嚇射撃をしたときには，81.6327秒以内にその場から離れることにしよう．


%%%%% 以下，空気抵抗の節（下書き） %%%%%
\section{空気抵抗を考えるとどうなるのか？理論編}

ここまでは空気抵抗がないと仮定してきた．

しかし現実世界には空気がある．空気の中を$400m/s$で飛ぶ弾丸は，当然ながら空気の抵抗を受ける．

ということは，ここまでの計算は少々弾丸に有利すぎたのではないか\footnote{弾丸に有利，とは？}？

というわけで空気抵抗を入れて計算し直してみよう．

\subsection{空気抵抗のモデル}

弾丸のように速く動く物体の場合，空気抵抗の大きさは速さの2乗に比例すると考えられる．
そこで空気抵抗の大きさを
\[
F=kv^2, \quad k=\frac{1}{2}\rho C_d A
\]
とおく．ここで$\rho$は空気の密度，$A$は弾丸の断面積，$C_d$は抗力係数と呼ばれる定数で，物体の形によって決まる\footnote{$\frac{1}{2}\rho v^2$は動圧と呼ばれ，速さ$v$で流れる空気の単位体積あたりの運動エネルギーにあたる．抗力係数$C_d$は，物体が受ける力が「動圧$\times$断面積」の何倍かを表す数として定義される．}．

空気抵抗はつねに運動と反対向きにはたらく．したがって上昇中は重力と空気抵抗がどちらも下向きに，落下中は重力が下向き，空気抵抗が上向きにはたらく．

以下，まず解析的な解を求める．その後で具体的な数値を与える．

\subsection{終端速度}

まずは簡単な落下中の弾丸を考える．落下の速さを$v$とおけば，下向きを正とした運動方程式は
\[
m\frac{dv}{dt}=mg-kv^2
\]
である．

速さ$v$が大きくなると空気抵抗$kv^2$も大きくなり，やがて重力$mg$とつりあう．つりあったところで加速度は0になり，弾丸はそれ以上速くならない．このときの速さを終端速度と呼び，$v_t$で表す．
\[
mg=kv_t^2 \iff v_t=\sqrt{\frac{mg}{k}}=\sqrt{\frac{2mg}{\rho C_d A}}
\]

つまりどんなに高いところから落としても，弾丸は$v_t$より速くは落ちてこない．

この$v_t$を使うと，上昇中と落下中の運動方程式はそれぞれ
\begin{align*}
\text{上昇中：}\quad \frac{dv}{dt}&=-g\left(1+\frac{v^2}{v_t^2}\right)\\
\text{落下中：}\quad \frac{dv}{dt}&=g\left(1-\frac{v^2}{v_t^2}\right)
\end{align*}
と書ける（上昇中は上向き，落下中は下向きをそれぞれ正とした）．

\subsection{最高高度}

知りたいのは「どの高さまで上がるか」なので，速さ$v$を時間$t$の関数としてではなく，高さ$y$の関数として考える．

速さは高さの時間変化だから$v=\dfrac{dy}{dt}$である．合成関数の微分法を使うと
\[
\frac{dv}{dt}=\frac{dv}{dy}\cdot\frac{dy}{dt}=v\frac{dv}{dy}
\]
となる．これを上昇中の運動方程式の左辺に代入すると
\[
v\frac{dv}{dy}=-g\left(1+\frac{v^2}{v_t^2}\right)
\]
を得る．

この式は，$v$を含む部分と$y$を含む部分を左右に分けることができる．このような形の微分方程式を変数分離形という．両辺を$1+v^2/v_t^2$で割り，$dy$を右辺に移すと
\[
\frac{v}{1+v^2/v_t^2}\,dv=-g\,dy
\]
となる\footnote{$dv$や$dy$を分数の分子・分母のように扱うのは形式的な操作である．正確には，両辺を$y$で積分してから，左辺に置換積分を使うと同じ結果になる．}．左辺の分母と分子に$v_t^2$を掛けて
\[
\frac{v_t^2\,v}{v_t^2+v^2}\,dv=-g\,dy
\]
と整理しておく．

次に両辺を積分する．発射した瞬間（$y=0$）の速さは$v_0$で，最高高度（$y=h$）では速さが0になる．したがって左辺は$v$について$v_0$から$0$まで，右辺は$y$について$0$から$h$まで積分する．
\[
\int_{v_0}^{0}\frac{v_t^2\,v}{v_t^2+v^2}\,dv=-g\int_0^h dy
\]

右辺は$-gh$である．

左辺は，分母$v_t^2+v^2$を$v$で微分すると$2v$になることに注目する．一般に
\[
\int\frac{f'(v)}{f(v)}\,dv=\log|f(v)|+C
\]
だから\footnote{$\log$は自然対数（底が$e$の対数）を表す．}，
\[
\int\frac{v}{v_t^2+v^2}\,dv=\frac{1}{2}\int\frac{2v}{v_t^2+v^2}\,dv=\frac{1}{2}\log(v_t^2+v^2)+C
\]
である（$v_t^2+v^2>0$なので絶対値は不要）．よって左辺は
\begin{align*}
\int_{v_0}^{0}\frac{v_t^2\,v}{v_t^2+v^2}\,dv
&=\frac{v_t^2}{2}\Big[\log(v_t^2+v^2)\Big]_{v_0}^{0}\\
&=\frac{v_t^2}{2}\left\{\log v_t^2-\log(v_t^2+v_0^2)\right\}\\
&=-\frac{v_t^2}{2}\log\frac{v_t^2+v_0^2}{v_t^2}\\
&=-\frac{v_t^2}{2}\log\left(1+\frac{v_0^2}{v_t^2}\right)
\end{align*}
となる．

左辺と右辺を等しいとおくと
\[
-\frac{v_t^2}{2}\log\left(1+\frac{v_0^2}{v_t^2}\right)=-gh
\]
だから，両辺を$-g$で割って，最高高度
\[
h=\frac{v_t^2}{2g}\log\left(1+\frac{v_0^2}{v_t^2}\right)
\]
を得る．
なんか一応解らしいものは出てきたが，これだと何なのかよく分からない．そこで次節では，落ちてくる弾の速さを求めてみよう．


\subsection{落ちてくるときの速さ}

次に，最高高度$h$で止まった弾丸が落ちてくる場合を考える．最高高度から下向きに測った落下距離を$s$とおく．落下を始める瞬間は$s=0$，地上に着く瞬間は$s=h$である．

落下中も$v=\dfrac{ds}{dt}$だから，上昇中と同じく合成関数の微分法で$\dfrac{dv}{dt}=v\dfrac{dv}{ds}$となり，落下中の運動方程式は
\[
v\frac{dv}{ds}=g\left(1-\frac{v^2}{v_t^2}\right)
\]
と書ける．これも変数分離形である．上昇中と同じ手順で変数を分けて整理すると
\[
\frac{v_t^2\,v}{v_t^2-v^2}\,dv=g\,ds
\]
となる．

落下を始める瞬間（$s=0$）の速さは0で，地上（$s=h$）に着いたときの速さを$v_f$とおく．左辺は$v$について$0$から$v_f$まで，右辺は$s$について$0$から$h$まで積分する．
\[
\int_{0}^{v_f}\frac{v_t^2\,v}{v_t^2-v^2}\,dv=g\int_0^h ds
\]

右辺は$gh$である．

左辺では，分母$v_t^2-v^2$を$v$で微分すると$-2v$になる．落下中の速さは終端速度より小さいので$v_t^2-v^2>0$であり，
\[
\int\frac{v}{v_t^2-v^2}\,dv=-\frac{1}{2}\int\frac{-2v}{v_t^2-v^2}\,dv=-\frac{1}{2}\log(v_t^2-v^2)+C
\]
となる．よって左辺は
\begin{align*}
\int_{0}^{v_f}\frac{v_t^2\,v}{v_t^2-v^2}\,dv
&=-\frac{v_t^2}{2}\Big[\log(v_t^2-v^2)\Big]_{0}^{v_f}\\
&=-\frac{v_t^2}{2}\left\{\log(v_t^2-v_f^2)-\log v_t^2\right\}\\
&=-\frac{v_t^2}{2}\log\left(1-\frac{v_f^2}{v_t^2}\right)
\end{align*}
である．

左辺と右辺を等しいとおくと
\[
-\frac{v_t^2}{2}\log\left(1-\frac{v_f^2}{v_t^2}\right)=gh
\]
となる．右辺の$h$に，前節で求めた最高高度を代入すると
\[
gh=\frac{v_t^2}{2}\log\left(1+\frac{v_0^2}{v_t^2}\right)
\]
だから，両辺を$-\dfrac{v_t^2}{2}$で割って
\[
\log\left(1-\frac{v_f^2}{v_t^2}\right)=-\log\left(1+\frac{v_0^2}{v_t^2}\right)=\log\frac{v_t^2}{v_t^2+v_0^2}
\]
を得る．両辺を$e$の指数にとると，$e^{\log x}=x$より
\[
1-\frac{v_f^2}{v_t^2}=\frac{v_t^2}{v_t^2+v_0^2}
\]
となる．
これを$v_f^2$について解くと
\[
\frac{v_f^2}{v_t^2}=1-\frac{v_t^2}{v_t^2+v_0^2}=\frac{v_0^2}{v_t^2+v_0^2}
\qquad\text{より}\qquad
v_f^2=\frac{v_0^2v_t^2}{v_0^2+v_t^2}
\]
となる．$v_f>0$だから，着地するときの速さは
\[
v_f=\frac{v_0v_t}{\sqrt{v_0^2+v_t^2}}
\]
である．

この式をよく見てほしい．分母と分子を$v_0$で割ると
\[
v_f=\frac{v_t}{\sqrt{1+v_t^2/v_0^2}}
\]
となる．分母の$\sqrt{1+v_t^2/v_0^2}$はつねに1より大きいので，$v_0$をどんなに大きくしても
\[
v_f<v_t
\]
だから，落ちてくる弾丸の速さは決して終端速度$v_t$を超えない．空気抵抗がないときには「撃った速さで戻ってくる」のだったが，空気抵抗があると「初速がどんなに大きくても$v_t$未満で戻ってくる」のである．

\section{具体的な数値を計算して実感する}


さて数式ばかり続いて読者も退屈しているだろうから，ここから楽しい具体的な数値計算に入る．

拳銃の弾として，よく使われる9mm弾を考えよう．ちなみに筆者は9mm弾を使うベレッタ92をグアムで撃ったことがある．撃つことよりも，弾倉を外したり装填したりするほうが楽しかった（撃つのは怖い）．
%要確認：弾丸の質量・抗力係数の値
質量を$m=8g=0.008kg$，直径を$9mm$とすれば，断面積は
\[
A=\pi\times(0.0045m)^2=6.36\times10^{-5}m^2
\]
である．空気の密度を$\rho=1.2kg/m^3$とし，抗力係数を$C_d=0.3$と仮定する\footnote{抗力係数は弾丸の形や，落下中に弾丸がどの向きを向いているかによって変わる．弾丸が横を向いて落ちてくれば$C_d$はもっと大きくなる．}．

すると終端速度は
\[
v_t=\sqrt{\frac{2\times0.008\times9.8}{1.2\times0.3\times6.36\times10^{-5}}}=82.7m/s
\]
となる．

初速$v_0=400m/s$を代入すると，最高高度は
\[
h=\frac{(82.7)^2}{2\times9.8}\log\left(1+\frac{400^2}{82.7^2}\right)=1115m
\]
落ちてきたときの速さは
\[
v_f=\frac{400\times82.7}{\sqrt{400^2+82.7^2}}=81.0m/s
\]
である．

空気抵抗がなければ$8163m$まで上昇するはずだった弾丸は，わずか$1115m$ほどしか上昇しない．そして$400m/s$で戻ってくるはずだった弾丸は，$81m/s$ほどで戻ってくる．

運動エネルギーで比べてみよう．発射した瞬間の運動エネルギーは
\[
K_0=\frac{1}{2}\times0.008\times400^2=640J
\]
であるのに対して，落ちてきた弾丸の運動エネルギーは
\[
K_f=\frac{1}{2}\times0.008\times81.0^2=26.3J
\]
である．発射した瞬間の約4\%にすぎない．

空気さん，ありがとう．

\subsection{それでも安心はできない}

では落ちてきた弾丸に当たっても平気かというと，そうとも言い切れない．

%要確認：皮膚を貫通する速さの値と出典
弾丸が人の皮膚を貫通するには，およそ$60m/s$程度の速さがあれば十分だと言われている．落ちてくる弾丸の速さ$81m/s$はこれを上回っている．
つまり当たりどころが悪ければ，落ちてきた弾丸でも十分に怪我をする．

%要確認：斜めに撃った場合の記述
また，ここでは真上に撃った場合だけを考えた．斜めに撃った弾丸は，最高点でも横向きの速度を保ったまま放物線に近い軌道で落ちてくるので，真上に撃った場合よりも速い速度で地上に戻ってくる可能性がある．

最後に，落ちてくるまでの時間も計算しておこう．運動方程式を時間について解くと，上昇にかかる時間$t_1$と落下にかかる時間$t_2$は
\[
t_1=\frac{v_t}{g}\arctan\frac{v_0}{v_t}, \qquad
t_2=\frac{v_t}{g}\,\mathrm{artanh}\frac{v_f}{v_t}
\]
となる．数値を代入すると$t_1=11.5$秒，$t_2=19.2$秒で，合計約$31$秒である．

空気抵抗がないときには約$82$秒だったから，実際には弾丸はずっと早く戻ってくる．

しかも上昇にかかる時間より，落下にかかる時間のほうが長い．空気抵抗がないときには両者が「ぴったり一致」していたのとは対照的である．

というわけで，結論を修正しなければならない．

これからはゾンビに威嚇射撃をしたときには，30秒以内にその場から離れることにしよう．


\section*{文献}
%要確認：出版社・刊行年
\bib 結城浩．2021．『数学ガールの物理ノート　ニュートン力学』SBクリエイティブ．



\end{document}
